\documentclass[10pt,a4paper]{amsart}
\usepackage[margin=19mm]{geometry}
\usepackage{amsmath,amssymb,mathtools}
\usepackage[scr=rsfs,cal=euler]{mathalfa}
\usepackage{xfrac}
\usepackage[unicode,hidelinks,bookmarks=false]{hyperref}
\newcommand{\ie}{i.e.\ }
\newcommand{\cf}{cf.\ }
\newcommand{\resp}{respectively\ }
\newcommand{\Subsection}{\subsection}
\providecommand{\nobreakdash}{\nobreak\hbox{-}}
\setlength{\emergencystretch}{2em}
\multlinegap0pt
\usepackage{color}
\usepackage{color}
\usepackage{txfonts}
\usepackage{enumitem}
\usepackage{bm}
 \newcommand{\red}[1]{{\color{red}#1}}
\newcommand{\R}{\mathbb{R}}
\newcommand{\C}{\mathbb{C}}
\newcommand{\Z}{\mathbb{Z}}
\newcommand{\N}{\mathbb{N}}
\newcommand{\ep}{\varepsilon}
\newcommand{\pa}{\partial}
\newcommand{\F}{\mathcal{F}}
\newcommand{\ee}{\mbox{\boldmath $1$}}
\DeclareMathOperator{\Arg}{Arg}
\DeclareMathOperator{\supp}{supp}
\DeclareMathOperator{\Id}{Id}
\newcommand{\lr}[1]{{}\langle{}#1{}\rangle{}}
\newcommand{\LR}[1]{{}\left\langle{}#1{}\right\rangle{}}
\newcommand{\ity}{\infty}
\newcommand{\dps}{\displaystyle}
\newcommand{\f}{\displaystyle\frac}
\setlength{\topmargin}{0mm}
\setlength{\oddsidemargin}{0mm}
\setlength{\evensidemargin}{0mm}
\setlength{\textwidth}{160mm}
\setlength{\textheight}{220mm}
\newtheorem{theorem}{Theorem}[section]
\newtheorem{lemma}[theorem]{Lemma}
\newtheorem{proposition}[theorem]{Proposition}
\newtheorem{corollary}[theorem]{Corollary}
\newtheorem{remark}{Remark}[section]
\newtheorem{notation}{Notation}
\newtheorem{definition}{Definition}[section]
\newtheorem{example}{Example}
\numberwithin{equation}{section}
\def\@cite#1#2{[{{\bfseries #1}\if@tempswa , #2\fi}]}
\begin{document}
\title[Global existence of solutions for nonlinear damped wave equa]{Global existence of solutions for nonlinear damped wave equations in an exterior domain with nonlinearities of derivative type}
\author{Tuan Anh Dao, Dinh Van Duong, Masahiro Ikeda}
\date{}
\maketitle
\noindent\emph{Source and licence.} Global existence of solutions for nonlinear damped wave equations in an exterior domain with nonlinearities of derivative type. Version arXiv:2608.04671v1 (CC BY 4.0). This material is distributed under the Creative Commons Attribution 4.0 International licence, http://creativecommons.org/licenses/by/4.0/, as recorded in the arXiv OAI licence metadata for this exact version and shown on the abstract page https://arxiv.org/abs/2608.04671v1. Full rights notice: licenses/arxiv-2608.04671-CC-BY-4.0.txt.\par\smallskip
\noindent\textit{Editorial scope.} Counted unit: the original \texttt{proposition} environment \texttt{prop:chain} at source lines 460--510 together with its complete proof at lines 512--831; the delivered document keeps source lines 226--832 so that every referenced label and every citation resolves inside it. Native editable source is the paper's own concatenated TeX; the AMS wrapper is editorial formatting only. No publisher class, upstream script or unrelated artwork is redistributed.
\section{Preliminaries}\label{Pre.Sec}

Throughout the paper, for $1\le q\le\infty$, we denote by $L^q(\Omega)$, the usual Lebesgue space with norm
$$
\|f\|_{L^q(\Omega)}
=
\begin{cases}
    \left(\displaystyle\int_\Omega |f(x)|^q,dx \right)^{1/q} &\text{ if } 1\le q<\infty, \\
    \operatorname*{ess,sup}_{x\in\Omega} |f(x)| &\text{ if } q= \infty.
\end{cases}
$$
We define the Dirichlet Laplacian through the theory of closed quadratic forms. Namely, let us consider the sesquilinear form
\begin{equation*}
%\label{eq:Dirichlet-form}
\mathfrak a[u,v]
:=
\int_{\Omega}
\nabla u(x)\cdot\overline{\nabla v(x)}\,dx,
\qquad
u,v\in D(\mathfrak a):=H_0^1(\Omega).
\end{equation*}
This form \(\mathfrak a\) is densely defined, symmetric, nonnegative, and closed on \(L^2(\Omega)\). By the representation theorem for closed nonnegative forms, there
exists a unique nonnegative self-adjoint operator  \(\mathcal{A}\) on \(L^2(\Omega)\) such that
\begin{equation*}
%\label{eq:form-representation}
\mathfrak a[u,v]
=
(\mathcal{A}u,v)_{L^2(\Omega)}
\end{equation*}
for every \(u\in D(\mathcal{A})\) and \(v\in H_0^1(\Omega)\), where
\begin{equation*}
%\label{eq:domain-A-form}
D(\mathcal{A})
=
\left\{
u\in H_0^1(\Omega):
\text{there exists } f\in L^2(\Omega)
\text{ such that }
\mathfrak a[u,v]=(f,v)_{L^2(\Omega)}
\text{ for all }v\in H_0^1(\Omega)
\right\},
\end{equation*}
and \(\mathcal{A}u:=f\). The uniqueness of \(f\) follows from the density of \(H_0^1(\Omega)\) in \(L^2(\Omega)\). This operator is the Dirichlet realization of the negative Laplacian, and we write $\mathcal{A}=-\Delta_{D,\Omega}$. In particular, in the distributional sense we have
\[
\mathcal{A}u=-\Delta u,
\qquad
u\in D(\mathcal{A}).
\]
If \(\Omega\) is of class \(\mathcal{C}^{1,1}\), the elliptic regularity yields
\begin{equation*}
%\label{eq:domain-A}
D(\mathcal{A})
=
H^2(\Omega)\cap H_0^1(\Omega),
\end{equation*}
with equivalence of the graph norm of \(\mathcal{A}\) and the \(H^2(\Omega)\)-norm on \(D(\mathcal{A})\).

Since \(\mathcal{A}\) is nonnegative and self-adjoint, the spectral theorem provides a projection-valued measure \(E_\mathcal{A}\) on \([0,\infty)\) such that
\begin{equation*}
%\label{eq:spectral-decomposition-A}
\mathcal{A}
=
\int_{[0,\infty)}
\lambda\,dE_\mathcal{A}(\lambda).
\end{equation*}
For \(u\in L^2(\Omega)\), define the finite positive Borel measure
\[
\mu_u(B)
:=
\big\|E_\mathcal{A}(B)u\big\|_{L^2(\Omega)}^2,
\qquad
B\subset[0,\infty)
\]
for every Borel set \(B\). For \(r\ge0\), the fractional power
\(\mathcal{A}^{r/2}\) is defined by
\begin{equation*}
%\label{eq:fractional-power}
\mathcal{A}^{r/2}u
:=
\int_{[0,\infty)}
\lambda^{r/2}\,dE_\mathcal{A}(\lambda)u
\end{equation*}
on the domain
\begin{equation*}
%\label{eq:fractional-domain}
D(\mathcal{A}^{r/2})
:=
\left\{
u\in L^2(\Omega):
\int_{[0,\infty)}
\lambda^r\,d\mu_u(\lambda)<\infty
\right\}.
\end{equation*}
Equivalently, it holds
\[
u\in D(\mathcal{A}^{r/2})
\quad\text{ if and only if }\quad
\int_{[0,\infty)}
\lambda^r
\,d\|E_\mathcal{A}(\lambda)u\|_{L^2(\Omega)}^2
<\infty
\]
and
\begin{equation*}
%\label{eq:fractional-power-norm}
\big\|\mathcal{A}^{r/2}u\big\|_{L^2(\Omega)}^2
=
\int_{[0,\infty)}
\lambda^r\,d\mu_u(\lambda)
\end{equation*}
for every \(u\in D(\mathcal{A}^{r/2})\). Thus, \(\mathcal{A}^{r/2}\) is a densely defined, nonnegative, self-adjoint, and closed operator on \(L^2(\Omega)\). For \(r\ge0\), we define
\begin{equation*}
%\label{eq:operator-Sobolev-space}
H^r(\mathcal{A})
:=
D(\mathcal{A}^{r/2})
\end{equation*}
and equip this space with the norm
\begin{equation*}
%\label{eq:operator-Sobolev-norm}
\|u\|_{H^r(\mathcal{A})}
:=
\|(\mathcal{I}+\mathcal{A})^{r/2}u\|_{L^2(\Omega)}.
\end{equation*}
By the spectral theorem, one sees that
\begin{equation*}
%\label{eq:operator-Sobolev-spectral-norm}
\|u\|_{H^r(\mathcal{A})}^2
=
\int_{[0,\infty)}
(1+\lambda)^r\,d\mu_u(\lambda),
\end{equation*}
together with the corresponding inner product
\begin{equation*}
%\label{eq:operator-Sobolev-inner-product}
(u,v)_{H^r(\mathcal{A})}
:=
\Big(
(\mathcal{I}+\mathcal{A})^{r/2}u,
(\mathcal{I}+\mathcal{A})^{r/2}v
\Big)_{L^2(\Omega)}.
\end{equation*}
Moreover, \(H^r(\mathcal{A})\) is a Hilbert space with respect to this inner product. 

For \(0<r<2\), \(D(\mathcal{A}^{r/2})\) is identified with the corresponding Dirichlet Sobolev space \(H^r_D(\Omega)\) and with the appropriate trace condition. Denoting the homogeneous H\"older space $\dot{\mathcal{C}}^\alpha(\mathbb R^2)$ for $0<\alpha<1$, whose seminorm is defined by
$$
\|f\|_{\dot{\mathcal{C}}^\alpha(\mathbb R^2)}
:=
\sup_{x\neq y}
\frac{|f(x)-f(y)|}
{|x-y|^\alpha}.
$$
We identify \(\mathbb C\) with \(\mathbb R^2\) and regard the nonlinear function \(\mathcal N\) as a mapping from \(\mathbb R^2\) into either
\(\mathbb R\) or \(\mathbb R^2\). Also, we introduce the following assumptions on the nonlinearity:

\begin{itemize}

\item $\mathcal N(0)=0$.

\item There exists a constant \(C_{\mathcal N}>0\) and \(p>1\) such that for all \(z_1,z_2\in\mathbb R^2\)
\begin{equation}
\label{eq:N-Lip}
|\mathcal N(z_1)-\mathcal N(z_2)|
\le
C_{\mathcal N}
\bigl(|z_1|+|z_2|\bigr)^{p-1}
|z_1-z_2|.
\end{equation}

\item Let \(1<\sigma<\min\{2,p\}\). We assume that $
\mathcal N\in \mathcal C^1(\mathbb R^2,\mathbb R^m)$ for $m\in\{1,2\},
$
and that there exists \(C_{\mathcal N}>0\) such that
\begin{equation}
\label{eq:DN-growth}
\big\|D\mathcal N(z)\big\|_{\mathcal C(\mathbb R^2,\mathbb R^m)}
\le
C_{\mathcal N}|z|^{p-1}
\end{equation}
for every \(z\in\mathbb R^2\). Moreover, \(D\mathcal N\) is assumed to be locally H\"older continuous
of order \(\sigma-1\) fulfilling
\begin{equation*}
%\label{eq:DN-local-holder}
\big\|D\mathcal N(z_1)-D\mathcal N(z_2)\big\|_{\mathcal C(\mathbb R^2,\mathbb R^m)}
\le
C_{\mathcal N}
\bigl(|z_1|+|z_2|\bigr)^{p-\sigma}
|z_1-z_2|^{\sigma-1}
\end{equation*}
for all \(z_1,z_2\in\mathbb R^2\).

\item If \(p>2\), we assume in addition that $\mathcal N\in \mathcal C^2(\mathbb R^2,\mathbb R^m)$
and for every \(z\in\mathbb R^2\) it holds
\begin{equation}
\label{eq:D2N-growth}
\big\|D^2\mathcal N(z)\big\|_{\mathcal C(\mathbb R^2,\mathbb R^m)}
\le
C_{\mathcal N}|z|^{p-2},
\end{equation}
where $D^2\mathcal N(z)\in
\mathcal C(\mathbb R^2,\mathbb R^m)$ denotes the second Fr\'echet derivative, regarded as a bounded bilinear mapping.
\end{itemize}
For later convenience, hereafter $C$ and $C_j$, with $j\in\N$, stand for suitable positive constants which may be changed from line to line.
%Moreover, for two given nonnegative functions $f$ and $g$, we write $f\lesssim g$ when $f\le Cg$, and $f \approx g$ when $g\lesssim f\lesssim g$.



\subsection{Fractional chain rule}

We first recall the fractional Leibniz rule on $\mathbb R^2$. This estimate will later be transferred to the setting of exterior domains via an extension argument.

\begin{lemma}[Fractional Leibniz rule on $\mathbb R^2$, see Lemma 2.7 in \cite{IkedaInuiWakasugi2026}]
\label{lem:Leibniz}
Let $\alpha\in (0,1)$. We assume
\[
f\in L^\infty(\mathbb R^2)\cap \dot{\mathcal{C}}^\alpha(\mathbb R^2)
\quad\text{ and }\quad
g\in \dot H^\alpha(\mathbb R^2)\cap L^2(\mathbb R^2),
\]
then there exists a positive constant $C$ depending only on $s$ such that
\begin{equation*}
%\label{eq:Leibniz}
\|fg\|_{\dot H^\alpha(\mathbb R^2)}
\le
C
\|f\|_{L^\infty(\mathbb R^2)}
\|g\|_{\dot H^\alpha(\mathbb R^2)}
+
C
\|f\|_{\dot{\mathcal{C}}^\alpha(\mathbb R^2)}
\|g\|_{L^2(\mathbb R^2)}.
\end{equation*}
\end{lemma}


\begin{proposition}[Fractional chain rule associated with the Dirichlet Laplacian]
\label{prop:chain}
Let $0<\sigma<2$ and $p>\sigma$. Assume that the nonlinear function $\mathcal{N}$ satisfies the above assumptions. Then, the following estimates hold:
\begin{itemize}
    \item[i)] If $f\in H^\sigma(\mathcal{A})\cap L^\infty(\Omega)$, then $\mathcal{N}[f]\in H^\sigma(\mathcal{A})$ and
    \begin{equation}
    \label{eq:chain}
    \big\|\mathcal{N}[f]\big\|_{H^\sigma(\mathcal{A})}
    \le
    C
    \|f\|_{L^\infty(\Omega)}^{p-1}
    \|f\|_{H^\sigma(\mathcal{A})}.
    \end{equation}
    \item[ii)] If $f,g\in D(\mathcal A^{\sigma/2})\cap L^\infty(\Omega)$, then $\mathcal N[f]-\mathcal N[g]\in D(\mathcal A^{\sigma/2})$
and
\begin{align}
&\bigl\|
\mathcal A^{\sigma/2}
\bigl(
\mathcal N[f]-\mathcal N[g]
\bigr)
\bigr\|_{L^2(\Omega)} \notag \\
&\quad\le
C
\left(
\|f\|_{L^\infty(\Omega)}^{p-1}
+
\|g\|_{L^\infty(\Omega)}^{p-1}
\right)
\bigl\|
\mathcal A^{\sigma/2}(f-g)
\bigr\|_{L^2(\Omega)}
\label{eq:Lip-chain-first}\\
&\qquad
+
C
\left(
\|f\|_{L^\infty(\Omega)}^{p-2}
+
\|g\|_{L^\infty(\Omega)}^{p-2}
\right)
\left(
\|\mathcal A^{\sigma/2}f\|_{L^2(\Omega)}
+
\|\mathcal A^{\sigma/2}g\|_{L^2(\Omega)}
\right)
\|v-w\|_{L^\infty(\Omega)}.
\notag
\end{align}
\end{itemize}
\end{proposition}

\begin{proof}
\textit{At first, let us verify the estimate \eqref{eq:chain}}. Namely, we take an extension operator \(\mathcal{E}\) satisfying $\mathcal{E}:H^\sigma_D(\Omega)\cap L^\infty(\Omega)
\to
H^\sigma(\mathbb R^2)\cap L^\infty(\mathbb R^2)$ and
\[
\big\|\mathcal{E}[f]\big\|_{H^\sigma(\mathbb R^2)}
\le C\|f\|_{H^\sigma_D(\Omega)},
\qquad
\big\|\mathcal{E}[f]\big\|_{L^\infty(\mathbb R^2)}
\le C\|f\|_{L^\infty(\Omega)}.
\]
Since $\|f\|_{H^\sigma(\mathcal{A})}
\sim
\|f\|_{H^\sigma(\Omega)}$, it suffices to prove
$$
\big\|\mathcal{N}[f]\big\|_{H^\sigma(\mathbb R^2)}
\le
C
\|f\|_{L^\infty(\mathbb R^2)}^{p-1}
\|f\|_{H^\sigma(\mathbb R^2)}.$$
For $0<\sigma<1$, using the Gagliardo characterization one has
$$
[f]_{\dot H^\sigma(\mathbb R^2)}^2
\sim
\iint_{\mathbb R^2\times\mathbb R^2}
\frac{|f(x)-f(y)|^2}{|x-y|^{2+2\sigma}}
\,dx\,dy.
$$
By \eqref{eq:N-Lip}, we obtain
$$
\big|\mathcal{N}[f(x)]-\mathcal{N}[f(y)]\big|
\le
C
\|f\|_{L^\infty(\mathbb R^2)}^{p-1}
|f(x)-f(y)|,
$$
which gives
$$
\big[\mathcal{N}[f]\big]_{\dot H^\sigma(\mathbb R^2)}
\le
C
\|f\|_{L^\infty(\mathbb R^2)}^{p-1}
[f]_{\dot H^\sigma(\mathbb R^2)}.
$$
Since $\mathcal{N}(0)=0$, it follows from \eqref{eq:N-Lip} that
$$
\big\|\mathcal{N}[f]\big\|_{L^2(\mathbb R^2)}
\le
C
\|f\|_{L^\infty(\mathbb R^2)}^{p-1}
\|f\|_{L^2(\mathbb R^2)},
$$
thus
$$
\big\|\mathcal{N}[f]\big\|_{H^\sigma(\mathbb R^2)}
\le
C
\|f\|_{L^\infty(\mathbb R^2)}^{p-1}
\|f\|_{H^\sigma(\mathbb R^2)}.
$$
Next, for $1<\sigma<2$ we put $\alpha=\sigma-1$. Since $\mathcal{N}$ is regarded as a $\mathcal{C}^1$ map from $\mathbb R^2$ to $\mathbb R^2$, we have
$$
\nabla \mathcal{N}[f] = D\mathcal{N}[f]\nabla f.
$$
By Lemma~\ref{lem:Leibniz}, we may estimate
\begin{align*}
\big\|\nabla \mathcal{N}[f]\big\|_{\dot H^\alpha(\mathbb R^2)}
&\le
C
\big\|D\mathcal{N}[f]\big\|_{L^\infty(\mathbb R^2)}
\|\nabla f\|_{\dot H^\alpha(\mathbb R^2)}
+
C
\big\|D\mathcal{N}[f]\big\|_{\dot{\mathcal{C}}^\alpha(\mathbb R^2)}
\|\nabla f\|_{L^2(\mathbb R^2)} \\
&\le
C
\|f\|_{L^\infty(\mathbb R^2)}^{p-1}
\|f\|_{H^\sigma(\mathbb R^2)}
\end{align*}
together with
$$
\big\|\mathcal{N}[f]\big\|_{L^2(\mathbb R^2)}
\le
C
\|f\|_{L^\infty(\mathbb R^2)}^{p-1}
\|f\|_{L^2(\mathbb R^2)},
$$
to arrive at
$$
\big\|\mathcal{N}[f]\big\|_{H^\sigma(\mathbb R^2)}
\le
C
\|f\|_{L^\infty(\mathbb R^2)}^{p-1}
\|f\|_{H^\sigma(\mathbb R^2)}.
$$
In this way, turning to $\Omega$ we may conclude the estimate \eqref{eq:chain} to complete our proof. \medskip

\textit{Next, let us verify the estimate \eqref{eq:Lip-chain-first}}. Indeed, we divide the proof into two steps as follows. \medskip

\noindent\textbf{$\bullet$ Step 1:} Setting $h:=f-g$ we rewrite the nonlinear difference in the form
\begin{align*}
\mathcal{N}[f(x)]-\mathcal{N}[g(x)]
=
\int_0^1
\frac{d}{d\theta}
\mathcal{N}[g(x)+\theta h(x)]
\,d\theta
&=
\int_0^1
D\mathcal{N}[g(x)+\theta h(x)]
h(x)\,d\theta \\
&= \int_0^1 D\mathcal{N}[h_\theta(x)]\, h(x)\,d\theta
\end{align*}
with $
h_\theta:=g+\theta h=(1-\theta)g+\theta f$. Then, applying the fractional Leibniz rule we obtain
\begin{equation}
\label{eq:product-B-h}
\big\|D\mathcal{N}[h_\theta] h\big\|_{\dot H^\sigma(\Omega)}
\le
C\big\|D\mathcal{N}[h_\theta]\big\|_{L^\infty(\Omega)}
\|h\|_{\dot H^\sigma(\Omega)}
+
C\|h\|_{L^\infty(\Omega)}
\big\|D\mathcal{N}[h_\theta]\big\|_{\dot H^\sigma(\Omega)}
\end{equation}
for \(1<\sigma<2\). Let us now estimate \(D\mathcal{N}[h_\theta]\) in \(L^\infty(\Omega)\) and \(\dot H^\sigma(\Omega)\). In particular, by \eqref{eq:DN-growth} one has
\[
\big|D\mathcal{N}[h_\theta(x)]\big|
\le
C|h_\theta(x)|^{p-1}
\le
C_p
\left(
|f(x)|^{p-1}+|g(x)|^{p-1}
\right),
\]
which leads to
\begin{equation}
\label{eq:Btheta-Linfty}
\big\|D\mathcal{N}[h_\theta]\big\|_{L^\infty(\Omega)}
\le
C
\left(
\|f\|_{L^\infty(\Omega)}^{p-1}
+
\|g\|_{L^\infty(\Omega)}^{p-1}
\right).
\end{equation}
Substituting \eqref{eq:Btheta-Linfty} into the first term of
\eqref{eq:product-B-h}, we obtain
\begin{align}
&\big\|D\mathcal{N}[h_\theta]\big\|_{L^\infty(\Omega)}
\|h\|_{\dot H^\sigma(\Omega)} \le
C
\left(
\|f\|_{L^\infty(\Omega)}^{p-1}
+
\|g\|_{L^\infty(\Omega)}^{p-1}
\right)
\|f-g\|_{\dot H^\sigma(\Omega)},
\label{eq:first-product-term}
\end{align}
which gives the first term on the right-hand side of
\eqref{eq:Lip-chain-first}. On the other hand, it holds from the standard Sobolev composition estimate that
\begin{equation}
\label{eq:composition-DN}
\big\|D\mathcal{N}[h_\theta]\big\|_{\dot H^\sigma(\Omega)}
\le
C
\big\|D^2\mathcal{N}[h_\theta]\big\|_{L^\infty(\Omega)}
\|h_\theta\|_{\dot H^\sigma(\Omega)}.
\end{equation}
By \eqref{eq:D2N-growth}, we derive
\[
\big|D^2\mathcal{N}[h_\theta(x)]\big|
\le
C|h_\theta(x)|^{p-2}
\le
C_p
\left(
|f(x)|^{p-2}+|g(x)|^{p-2}
\right)
\]
since \(p\ge2\). Thus, it follows
\begin{equation*}
%\label{eq:D2N-ztheta}
\big\|D^2\mathcal N[h_\theta]\big\|_{L^\infty(\Omega)}
\le
C
\left(
\|f\|_{L^\infty(\Omega)}^{p-2}
+
\|g\|_{L^\infty(\Omega)}^{p-2}
\right),
\end{equation*}
which implies from \eqref{eq:composition-DN} that
\begin{align*}
\big\|D\mathcal{N}[h_\theta]\big\|_{\dot H^\sigma(\Omega)}
&\le
C
\left(
\|f\|_{L^\infty(\Omega)}^{p-2}
+
\|g\|_{L^\infty(\Omega)}^{p-2}
\right) \left(
\|f\|_{\dot H^\sigma(\Omega)}
+
\|g\|_{\dot H^\sigma(\Omega)}
\right)
%\label{eq:Btheta-Hsigma}
\end{align*}
thanks to the triangle inequality $\|h_\theta\|_{\dot H^\sigma(\Omega)} \le
\|f\|_{\dot H^\sigma(\Omega)}
+
\|g\|_{\dot H^\sigma(\Omega)}$. Therefore, the second term of \eqref{eq:product-B-h} is controlled by
\begin{align}
&\|h\|_{L^\infty(\Omega)}
\big\|D\mathcal{N}[h_\theta]\big\|_{\dot H^\sigma(\Omega)}
\le
C
\left(
\|f\|_{L^\infty(\Omega)}^{p-2}
+
\|g\|_{L^\infty(\Omega)}^{p-2}
\right)
\left(
\|f\|_{\dot H^\sigma(\Omega)}
+
\|g\|_{\dot H^\sigma(\Omega)}
\right)
\|f-g\|_{L^\infty(\Omega)}.
\label{eq:second-product-term}
\end{align}
This is exactly the second term required in
\eqref{eq:Lip-chain-first}. \medskip

\noindent\textbf{$\bullet$ Step 2:} Combining \eqref{eq:product-B-h},
\eqref{eq:first-product-term} and
\eqref{eq:second-product-term}, we obtain the following uniform estimate for \(\theta\in[0,1]\):
\begin{align*}
\big\|D\mathcal{N}[h_\theta] h\big\|_{\dot H^\sigma(\Omega)}
&\le
C
\left(
\|f\|_{L^\infty(\Omega)}^{p-1}
+
\|g\|_{L^\infty(\Omega)}^{p-1}
\right)
\|f-g\|_{\dot H^\sigma(\Omega)}
\notag\\
&\quad+
C
\left(
\|f\|_{L^\infty(\Omega)}^{p-2}
+
\|g\|_{L^\infty(\Omega)}^{p-2}
\right) \Big(
\|f\|_{\dot H^\sigma(\Omega)}
+
\|g\|_{\dot H^\sigma(\Omega)}
\Big)
\|f-g\|_{L^\infty(\Omega)}.
%\label{eq:Btheta-h-final}
\end{align*}
Applying Minkowski's integral inequality we arrive at
\begin{align}
\big\|\mathcal N[f]-\mathcal N[g]\big\|_{\dot H^\sigma(\Omega)}
&=
\left\|
\int_0^1 D\mathcal{N}[h_\theta] h\,d\theta
\right\|_{\dot H^\sigma(\Omega)}
\le
\int_0^1
\big\|D\mathcal{N}[h_\theta] h\big\|_{\dot H^\sigma(\Omega)}
\,d\theta \notag \\
&\le
C
\left(
\|f\|_{L^\infty(\Omega)}^{p-1}
+
\|g\|_{L^\infty(\Omega)}^{p-1}
\right)
\|f-g\|_{\dot H^\sigma(\Omega)} \notag \\
&\qquad+
C
\left(
\|f\|_{L^\infty(\Omega)}^{p-2}
+
\|g\|_{L^\infty(\Omega)}^{p-2}
\right) \left(
\|f\|_{\dot H^\sigma(\Omega)}
+
\|g\|_{\dot H^\sigma(\Omega)}
\right)
\|f-g\|_{L^\infty(\Omega)}.
\label{eq:whole-Sobolev-difference}
\end{align}
Finally, let us turn to the Dirichlet Laplacian to get the equivalent relation
\begin{equation*}
%\label{eq:A-Sobolev-equivalence-difference}
\|f\|_{\dot H_D^\sigma(\Omega)}
\sim
\|\mathcal A^{\sigma/2}f\|_{L^2(\Omega)}
\end{equation*}
for \(1<\sigma<2\). Consequently, one obtains the following estimates:
\begin{align*}
    \|f-g\|_{\dot H_D^\sigma(\Omega)}
&\lesssim
\big\|\mathcal A^{\sigma/2}(f-g)\big\|_{L^2(\Omega)}, \\
\|f\|_{\dot H_D^\sigma(\Omega)}
+
\|g\|_{\dot H_D^\sigma(\Omega)}
&\lesssim
\big\|\mathcal A^{\sigma/2}f\big\|_{L^2(\Omega)}
+
\big\|\mathcal A^{\sigma/2}g\big\|_{L^2(\Omega)}.
\end{align*}
In this way, plugging the two previous estimates into \eqref{eq:whole-Sobolev-difference} we may conclude the estimate \eqref{eq:Lip-chain-first}. Hence, our proof is established.
\end{proof}


\begin{thebibliography}{99}
\bibitem[IkedaInuiWakasugi2026]{IkedaInuiWakasugi2026} M. Ikeda, T. Inui, Y. Wakasugi, Existence of solutions to the semilinear damped wave equation with non-\(L^2\) slowly decaying data: polynomial nonlinearity case, \textit{J. Evol. Equ.}, \textbf{26} (2026), Article 37.
\end{thebibliography}
\end{document}
